\section{Medical Report OCR \& Structured Data Extraction}
\label{sec:ocr_pipeline}

To address the challenge of integrating paper-based medical records and scanned clinical laboratory reports into our digital system, we developed a high-precision \ac{OCR}, LLM-based data extraction pipeline. This system leverages vision-language models (\ac{VLM}s) to handle complex document structures, such as nested tables and specific medical units, converting unstructured visual data into structured \ac{JSON} and \ac{HL7} \ac{FHIR} R4 resources.

\subsection{System Architecture and Inference Backend}
\label{subsec:ocr_backend}

The pipeline follows a decoupled architecture consisting of a high-performance inference backend and a Python-based processing frontend. The core intelligence is powered by the \texttt{PaddleOCR-VL-1.5-0.9B} model, served via an official \ac{vLLM} container. 

\begin{itemize}
    \item \textbf{Model Selection:} The \texttt{PaddleOCR-VL-1.5-0.9B} model was chosen for its optimization in document understanding and joint layout analysis.
    \item \textbf{Hardware Optimization:} To ensure deployment feasibility on consumer-grade hardware (8GB \ac{VRAM} \ac{GPU}s), several optimization techniques are employed:
    \begin{itemize}
        \item \textbf{Memory Management:} \ac{GPU} memory utilization is capped at 70\% to prevent Out-of-Memory (\ac{OOM}) errors during peak inference.
        \item \textbf{Eager Execution:} \texttt{enforce\_eager=True} is used to reduce overhead on lower-end architectures.
        \item \textbf{Token Constraints:} The maximum model length is restricted to 4096 tokens to balance the context window against the memory footprint.
    \end{itemize}
\end{itemize}

\subsection{OCR Processing Pipeline}
\label{subsec:ocr_processing}

The processing pipeline handles document ingestion, preprocessing, and semantic extraction through the following stages:

\begin{enumerate}
    % \item \textbf{Preprocessing:} \ac{PDF} documents are rendered at 110 \ac{DPI} to maintain legibility while keeping visual token counts manageable. Images are automatically downscaled (typically between 960px and 2048px) to fit within the \ac{VLM}'s visual context window.
    \item \textbf{Core Extraction:} The system performs joint layout analysis and text recognition to identify hierarchical table structures and medical units (e.g., mg/dL, mmol/L).
    \item \textbf{Output Format:} The pipeline generates:
    \begin{itemize}
        % \item \textbf{Structured JSON:} Containing raw extraction data, bounding boxes, and confidence scores.
        \item \textbf{Markdown:} A human-readable representation of the document structure including reconstructed tables.
        % \item \textbf{Visual Verification:} Images with overlaid bounding boxes for manual auditing.
    \end{itemize}
\end{enumerate}

\section{OCR Clinical Lab Report output to HL7 FHIR R4 Pipeline}
\label{sec:fhir_pipeline}

Extracting raw text via \ac{OCR} is insufficient for clinical interoperability. We implemented a two-stage hybrid pipeline to transform unstructured summaries into highly structured \ac{HL7} \ac{FHIR} R4 transaction Bundles. This approach combines the natural language understanding (\ac{NLU}) of a localized \ac{LLM} (\texttt{Qwen2.5-3B-Instruct}) with deterministic Python logic.

\subsection{The Multi-Stage Transformation Strategy}
\label{subsec:two_stage_fhir}

\begin{figure}[H]
    \centering
    \includegraphics[width=1\linewidth]{AI/OCR-FHIR_Mapper.png.jpeg}
    % \caption{Interface of the MedMCP system demonstrating a Local RAG (Retrieval-Augmented Generation) response. The output includes synthesized clinical guidelines supported by specific evidence sources and PMIDs.}
    \label{fig:OCR-FHIR_Mapper}
\end{figure}

Relying solely on \ac{LLM}s for \ac{FHIR} generation often results in "hallucinated" clinical codes (LOINC, SNOMED) or syntactically invalid \ac{JSON}. Our architecture mitigates this by separating intent from assembly:

\begin{itemize}
    \item \textbf{Stage 1: Guided LLM Extraction:} The \texttt{Qwen2.5-3B} model is restricted strictly to parsing the \ac{OCR} text into a flat, intermediate \ac{JSON} schema. This stage uses explicit schema enforcement and an adaptive retry mechanism with increasing temperature if parsing fails.
    \item \textbf{Stage 2: Deterministic FHIR Builder:} A rule-based Python layer handles the mapping of extracted test names to authoritative \ac{LOINC} codes and translates units to \ac{UCUM} standards. This stage eliminates the risk of generative hallucination in medical coding.
\end{itemize}

\subsection{Terminology Mapping and Validation}
\label{subsec:fhir_mapping}

The deterministic builder implements several medical-grade validation logic:
\begin{enumerate}
    \item \textbf{LOINC Mapping:} Test names (e.g., "Total Leukocyte Count") are mapped via a hardcoded dictionary to authoritative codes (e.g., 6690-2).
    \item \textbf{Unit Scaling:} Values are mathematically scaled to \ac{UCUM} standards (e.g., scaling "lakhs/cumm" by 100 to map to standard units).
    % \item \textbf{Interpretation Cross-Validation:} The system extracts "High/Low" flags from the report but also re-calculates them using numeric values and reference ranges. Any disagreement triggers a warning, with the calculated value acting as the source of truth.
\end{enumerate}

\subsection{Bundle Assembly and Strict Audit}
\label{subsec:fhir_audit}

The final output is a \ac{FHIR} transaction Bundle containing \texttt{Patient}, \texttt{DiagnosticReport}, \texttt{Organization}, \texttt{Practitioner}, and multiple \texttt{Observation} resources. Before submission, the bundle undergoes two layers of validation:
\begin{itemize}
    \item \textbf{Structural Validation:} Using \texttt{fhir.resources} (Pydantic v2) to ensure adherence to the \ac{HL7} \ac{FHIR} R4 schema.
    \item \textbf{Semantic Validation:} Custom business rules ensure clinical integrity (e.g., no dangling references, mandatory \ac{LOINC} codes, and secure mapping of patient gender).
\end{itemize}
